\chapter{Auxiliary Development}

While not directly in the scope of our project, there was a lot
of "background" work that had to be done, so as to achieve the
goals we had laid out for our project. This auxiliary development
was made across different open source projects, to the which we
have been able to make contributions.

\section{Kamaillio Project}

The Kamaillio Project, previously known as OpenSER is a Session Initiation
Protocol (SIP). It is of great importance to our project namely due to its
telephony and IMS features. These capabilities are what allow for the
development of the SMSC that we deployed to IT's infrastructure, which was
nonexistent. To put it simply, one of the auxiliary developments that we had to
do in our project was the implementation of SMS throughout IT's network.

Our contribution can be visible in PR
\href{https://github.com/kamailio/kamailio/pull/4165}{\#4165}.

\section{The NEF Emulator Project}

This project was always intended to be a part of the final deployed stack,
mainly because of the location aspect that needs to be emulated, since IT does
not have a Location Management Function (LMF) in the 5G Cores, meaning we do
not have access to real location information. However, as soon as we started
development of the location APIs in our gateway we quickly found out that a lot
of issues needed to be addressed, and no collaborator from the project would
fix them for us, thus our participation in the development efforts. Initially
the whole location part of the NEF Emulator had to be updated to 5G Release 18,
since it was following outdated standards that were not compatible with our
needs. Then, a new endpoint was added to allow for the polling of a given
device's location. It was also needed to completely redo the logic for the
location subscriptions, since the existing codebase had a synchronous, blocking
threaded approach for both the simulation of the movement of the UEs and the
handling of the subscriptions. This meant that if either a subscription
notification or a missed update on a movement simulation failed, the whole
program would stop working. Another problem was that, due to the way it was
implemented, only one subscription was possible per device, while the standards
and our use case needed more than one subscription per device.

All the development on location can be seen on PRs
\href{https://github.com/ATNoG/NEF_emulator/pull/4}{\#4},
\href{https://github.com/ATNoG/NEF_emulator/pull/5}{\#5},
\href{https://github.com/ATNoG/NEF_emulator/pull/13}{\#13},
\href{https://github.com/ATNoG/NEF_emulator/pull/14}{\#14} and
\href{https://github.com/ATNoG/NEF_emulator/pull/15}{\#15}.

Following the issues we had with the location side of the
project, it was decided that the NEF Emulator would also be
mandatory to implement the Communication Quality and Device
Reachability APIs. To achieve this, it was also necessary to
update these functions to 5G's Release 18, since once again
these were following an outdated standard that was not compatible
with our use case, as seen in PR
\href{https://github.com/ATNoG/NEF_emulator/pull/8}{\#8}.

In addition to updating the schemas, a small fix had to be made
in an endpoint, in PR
\href{https://github.com/ATNoG/NEF_emulator/pull/10}{\#10}
and support for the \textbf{ROAMING\_STATUS} monitoring type,
which added a visiting PLMN ID to a given device, which allowed
for the recognition of whether the device was in roaming status
or not. This is visible in PR
\href{https://github.com/ATNoG/NEF_emulator/pull/16}{\#16}.


\section{IT Infrastructure Integration}

While it was needed to contribute to the different projects in order for them
to support the features that we needed, it was even more important to integrate
these with IT's infrastructure. Firstly, we had to configure, package, and
deploy the IMS and SMSC. Then the 5G Network was configured to provide a Data
Network Name (DNN) according to a GSMA Norm \cite{NG113}, and lastly the
Session Management Function (SMF) was configured to announce the proper address
\cite{TS23501} \cite{TS23228}. This implied not only connecting the SMSC and
IMS to the network, but also configuring the network to announce the correct
parameters for allowing the devices.

Then, we had to integrate the newly developed QoS features from
NEF to IT's 5G Huawei Core. This ended up proving to be a
challenge, specially since the only interaction allowed with the
core is through a Web UI. While IT has developed a custom API to
interact with this Web UI, using Selenium, this proved to be very
unreliable and caused several issues. The problem was exacerbated
by the lack of documentation, which severely hindered the speed of
implementation.

Finally, to perform our final demo, we integrated our project with IT's CCAM
Stack, using the developed features of digital twin in the NEF Emulator. IT's
CCAM\footnote{\url{https://ccam.av.it.pt/about}} is a tool that enables
vehicles to interact not only with each other but also with the surrounding
infrastructure, this allows for a seamless exchange of data. In order to send
Maneuver Coordination Messages \cite{MCM} and Collective Perception Messages
(CPM), an MQTT Broker is used. These messages contain the location of the
vehicle, which needs to be imported into NEF Emulator. Therefore, we have
deployed a pod that consumes these messages from the MQTT broker and injects
the relevant information into NEF Emulator. This way, we have access to the
location of the vehicle through the NEF interface, just like we would have in a
real 5G network deployment.
